\documentclass{projectUESPI}
\usepackage[utf8]{inputenc}
\usepackage[T1]{fontenc}
\usepackage{geometry}
\usepackage{graphicx}
\usepackage{float}
\usepackage{hyperref}
\usepackage[english,indonesian]{babel}

\title{Optimasi Pemrosesan Template LaTeX Menggunakan Pendekatan Berbasis AI}
\author{Academic Support Agent}
\date{\today}

\begin{document}
\maketitle

\begin{abstract}

\end{abstract}

\chapter{PENDAHULUAN}


\chapter{TINJAUAN PUSTAKA}


% Original content (commented out):
% \documentclass{projectUESPI}
\usepackage[T1]{fontenc}
\usepackage[utf8]{inputenc}
\usepackage[portuguese]{babel}
\usepackage{tabularx}
\usepackage{multirow}
\usepackage{float}
\usepackage{graphics}
\usepackage{subfigure}
\usepackage{graphicx}
\usepackage{float}
\usepackage{hyperref}

% \usepackage[alf]{abntex2cite}
\usepackage[alf,
versalete,
abnt-emphasize = bf, % destaca o titulo da revista ou livro em negrito;
abnt-etal-list = 3, % trabalhos com mais de 3 autores recebem et al.,;
abnt-etal-text = it, % escreve o et al., em italico;
%abnt-and-type = &, % usa o carater '&' no lugar de 'e' para mais de um autor;
abnt-doi=link,
abnt-title-command=no,
abnt-last-names = abnt, % trata sobrenomes 'estritamente' conforme a ABNT; e
abnt-repeated-author-omit = no % autores com + de uma entrada recebem '____.'
]{abntex2cite}
\usepackage[brazilian,hyperpageref]{backref}
% \renewcommand{\ABNTbibliographyname}{Referências2}
% \renewcommand{\ABNTbibliographyname}{Referências}



% \renewcommand{\bibtextitlecommand}{} %remove no nome referências antes da bibliografia.

% referência util: http://mirrors.ibiblio.org/CTAN/macros/latex/contrib/abntex2/doc/abntex2cite-alf.pdf

\addto\captionsportuguese{\renewcommand{\refname}{Referências Bibliográficas }}




\campus{Campus Dra. Josefina Demes}
\title{\color{red}TÍTULO DO PROJETO DE DESENVOLVIMENTO TECNOLÓGICO\color{black}}
\author{\color{red}Autores\color{black}}
\cidade{Floriano/PI}
\mesdate{\color{red}Mês Ano\color{black}}
% \usepackage{helvet}
% \renewcommand{\familydefault}{\sfdefault}
\begin{document}
\maketitle
\newpage
\tableofcontents
\thispagestyle{empty}
\newpage
\setcounter{page}{1}

\section{Apresentação e Justificativa}
\input{sections/apresentacao-justificativa}
\newpage

\section{Problemas e Soluções}
\input{sections/problemas-solucoes}
\newpage 


\section{Objetivos}
\input{sections/objetivos}
\newpage

\section{Métodos}
\input{sections/metodos-e-procedimentos}
\newpage

\section{Produtos e/ou Resultados Esperados}
\input{sections/resultados-esperados}
\newpage

\section{Cronograma de Execução do Projeto de Desenvolvimento}
\input{sections/cronograma}
\newpage

\addcontentsline{toc}{section}{Referências Bibliográficas}
\bibliography{referencias.bib}

\newpage

\end{document}


\end{document}